% =====================================================================
\chapter{Regla de la cadena y modelos de Markov}
\label{tema:2}
% =====================================================================

Este tema descompone la probabilidad de una secuencia en factores de la forma $P(x_t\mid x_{1:t-1})$ y presenta la aproximación de Markov que da lugar a los modelos de $n$-gramas. Los resultados de probabilidad siguen a~\cite[§1.3]{bt}; la parte de $n$-gramas, a~\cite[§3.1]{jm3}.

\section{Probabilidad condicionada}

\begin{definicion}[Probabilidad condicionada]
Sea $(\Omega,\mathcal{F},P)$ un espacio de probabilidad y sean $A,B\in\mathcal{F}$ con $P(A)>0$. La \emph{probabilidad de $B$ condicionada a $A$} es
\[
P(B\mid A)=\frac{P(A\cap B)}{P(A)}.
\]
\end{definicion}

Para $A$ fijo, la aplicación $B\mapsto P(B\mid A)$ es a su vez una medida de probabilidad sobre $(\Omega,\mathcal{F})$~\cite[§1.3]{bt}. Despejando de la definición se obtiene la \emph{regla del producto}, $P(A\cap B)=P(A)\,P(B\mid A)$.

\section{La regla de la cadena}

\begin{teorema}[Regla de la cadena]\label{teo:cadena}
Sean $A_1,\dots,A_T$ sucesos con $P(A_1\cap\dots\cap A_{T-1})>0$. Entonces
\[
P\Bigl(\,\bigcap_{t=1}^{T}A_t\Bigr)=\prod_{t=1}^{T}P\bigl(A_t\mid A_1\cap\dots\cap A_{t-1}\bigr),
\]
donde el factor con $t=1$ se interpreta como $P(A_1)$.
\end{teorema}

\begin{proof}
Por inducción sobre $T$. Para $T=1$ ambos miembros valen $P(A_1)$. Supóngase cierto para $T-1$ y sea $B=A_1\cap\dots\cap A_{T-1}$. Como $P(B)>0$, la regla del producto da
$P(B\cap A_T)=P(B)\,P(A_T\mid B)$. Por la hipótesis de inducción (aplicable porque $P(A_1\cap\dots\cap A_{t-1})\ge P(B)>0$ para todo $t\le T-1$), $P(B)=\prod_{t=1}^{T-1}P(A_t\mid A_1\cap\dots\cap A_{t-1})$. Sustituyendo se obtiene el enunciado.
\end{proof}

\begin{corolario}\label{cor:cadena-secuencias}
Para un modelo de lenguaje y una secuencia $x_{1:T}$ con $P(x_{1:T-1})>0$,
\[
P(x_{1:T})=\prod_{t=1}^{T}P(x_t\mid x_{1:t-1}).
\]
\end{corolario}

\begin{proof}
Basta aplicar el \cref{teo:cadena} a los sucesos $A_t=\{\text{el token en la posición } t \text{ es } x_t\}$.
\end{proof}

La igualdad del \cref{cor:cadena-secuencias} es exacta. Todo modelo de lenguaje autorregresivo, desde un modelo de bigramas hasta un Transformer, consiste en una forma concreta de calcular cada factor $P(x_t\mid x_{1:t-1})$.

\section{Hipótesis de Markov y modelos de \texorpdfstring{$n$}{n}-gramas}

El contexto $x_{1:t-1}$ crece con $t$, y el número de contextos posibles crece exponencialmente (\cref{prop:contar-secuencias}). La solución clásica consiste en truncar el contexto~\cite[§3.1.1]{jm3}.

\begin{definicion}[Modelo de $n$-gramas]\label{def:ngramas}
Sea $n\ge 1$. Un modelo de lenguaje satisface la \emph{hipótesis de Markov de orden $n-1$} si
\[
P(x_t\mid x_{1:t-1})=P(x_t\mid x_{t-n+1:t-1})\qquad\text{para todo } t,
\]
es decir, si la distribución del token siguiente depende solo de los $n-1$ tokens anteriores. El modelo se llama de \emph{unigramas} ($n=1$), de \emph{bigramas} ($n=2$) o de \emph{trigramas} ($n=3$).
\end{definicion}

\begin{definicion}[Matriz estocástica]
Una matriz $P\in\R^{K\times K}$ es \emph{estocástica por filas} si $P_{ij}\ge 0$ para todo $i,j$ y $\sum_{j}P_{ij}=1$ para todo $i$. Un modelo de bigramas queda determinado por la matriz estocástica $P_{ij}=P(x_t=j\mid x_{t-1}=i)$, llamada \emph{matriz de transición}.
\end{definicion}

\begin{proposicion}\label{prop:parametros}
Un modelo de $n$-gramas sobre un vocabulario de tamaño $K$ se especifica con una tabla de $K^{n}$ probabilidades, de las cuales $K^{n-1}(K-1)$ son libres.
\end{proposicion}

\begin{proof}
Hay $K^{n-1}$ contextos posibles y, para cada uno, una distribución sobre $K$ tokens: $K^{n-1}\cdot K=K^{n}$ entradas. Cada distribución está sujeta a una restricción de suma igual a~1, que fija una entrada a partir de las demás; quedan $K-1$ libres por contexto.
\end{proof}

Con $K=5\cdot 10^{4}$ y $n=3$, la tabla tendría $1{,}25\cdot 10^{14}$ entradas. Esta explosión combinatoria es la razón por la que los modelos modernos sustituyen la tabla por una función con parámetros compartidos (módulos~9 y~13).

\begin{ejemplo}[Un modelo de bigramas]\label{ej:bigramas}
Con el vocabulario del \cref{ej:luna} y el corpus «ana», «ala», «luna», el estimador de máxima verosimilitud del \cref{tema:3} da la matriz de transición de la \cref{fig:cadena}. Por el \cref{cor:cadena-secuencias},
\[
P(\text{ana})=P(\texttt{a}\mid\fin)\,P(\texttt{n}\mid\texttt{a})\,P(\texttt{a}\mid\texttt{n})\,P(\fin\mid\texttt{a})=\tfrac{2}{3}\cdot\tfrac{1}{5}\cdot 1\cdot\tfrac{3}{5}=\tfrac{2}{25}=0{,}08.
\]
Del mismo modo, $P(\text{luna})=0{,}10$, $P(\text{ala})=0{,}04$ y $P(\text{lana})=0{,}02$. La palabra «lana» no aparece en el corpus y recibe probabilidad positiva: el modelo combina transiciones observadas.
\end{ejemplo}

\begin{figure}[h]
\centering
\begin{tikzpicture}[->,>={Stealth[length=2mm]},node distance=2.4cm,every state/.style={minimum size=9mm,draw=black,thick},auto,font=\small]
  \node[state] (f) {\fin};
  \node[state] (a) [right=of f] {\texttt{a}};
  \node[state] (n) [above right=1.1cm and 2.0cm of a] {\texttt{n}};
  \node[state] (l) [below right=1.1cm and 2.0cm of a] {\texttt{l}};
  \node[state] (u) [right=3.2cm of a] {\texttt{u}};
  \path (f) edge[bend left=18] node {$2/3$} (a)
        (f) edge[bend right=40] node[below] {$1/3$} (l)
        (a) edge[bend left=18] node {$3/5$} (f)
        (a) edge[bend left=10] node {$1/5$} (n)
        (a) edge[bend left=10] node {$1/5$} (l)
        (n) edge[bend left=10] node {$1$} (a)
        (l) edge[bend left=10] node {$1/2$} (a)
        (l) edge node[right] {$1/2$} (u)
        (u) edge node[right] {$1$} (n);
\end{tikzpicture}
\caption{Diagrama de transición del modelo de bigramas del \cref{ej:bigramas}. Cada arco $i\to j$ lleva $P(j\mid i)$; las probabilidades que salen de cada estado suman~1.}
\label{fig:cadena}
\end{figure}

\section{El símbolo de fin y la normalización}

La \cref{def:ml} exige que las probabilidades de \emph{todas} las secuencias finitas sumen~1. El símbolo \fin\ de fin de secuencia es imprescindible para ello~\cite[§3.1, ejercicio 3.5]{jm3}.

\begin{proposicion}\label{prop:longitud-fija}
Sea un modelo de bigramas sin símbolo de fin, con distribución inicial $\pi$ sobre $\V$ y matriz de transición estocástica $P$. Para cada longitud $L\ge 1$,
\[
\sum_{x_{1:L}\in\V^{L}}\pi(x_1)\prod_{t=2}^{L}P_{x_{t-1}x_t}=1.
\]
En consecuencia, la suma sobre todas las longitudes diverge y el modelo no define una distribución sobre $\V^{*}$.
\end{proposicion}

\begin{proof}
Por inducción sobre $L$. Para $L=1$ la suma es $\sum_{x_1}\pi(x_1)=1$. Para pasar de $L$ a $L+1$, se suma primero sobre $x_{L+1}$: el factor nuevo es $\sum_{x_{L+1}}P_{x_Lx_{L+1}}=1$ porque $P$ es estocástica, y queda la suma de longitud $L$, igual a~1. Sumando sobre $L=1,2,\dots$ se obtiene $1+1+\dots=\infty$.
\end{proof}

\begin{observacion}
Con el símbolo de fin, la generación se detiene la primera vez que aparece \fin, y cada longitud recibe solo una parte de la masa total. Si desde cualquier token se alcanza \fin\ con probabilidad positiva, la masa que queda en secuencias de longitud $L$ decrece geométricamente con $L$ y la suma total es~1; el argumento completo es el de las cadenas de Markov absorbentes~\cite[cap.~7]{bt}.
\end{observacion}

\begin{lecturas}
\item Bertsekas y Tsitsiklis~\cite{bt}, §1.3 «Conditional Probability», regla de la multiplicación.
\item MIT 6.041SC~\cite{mit6041}, unidad~I, clase~2 (probabilidad condicionada).
\item Jurafsky y Martin~\cite{jm3}, §3.1 «N-Grams» y §3.1.1 «The Markov assumption».
\item Zhang \emph{et al.}~\cite{d2l}, §9.3.1.1 «Markov Models and $n$-grams».
\end{lecturas}

\section*{Ejercicios del tema 2}
\addcontentsline{toc}{section}{Ejercicios del tema 2}

\begin{ejercicio}{1}{eje:2-casillas}
Con un vocabulario de 27 tokens (26 letras y \fin), ¿cuántas entradas tiene la tabla de un modelo de bigramas? ¿Y la de uno de trigramas? ¿Cuántas de ellas son parámetros libres en cada caso?
\end{ejercicio}

\begin{ejercicio}{1}{eje:2-vf}
Justifique si son verdaderas o falsas: (a) la regla de la cadena es una aproximación; (b) cada fila de una matriz de transición suma~1.
\end{ejercicio}

\begin{ejercicio}{2}{eje:2-sol-matriz}
Corpus: «sol», «sal», «los». Con $\V=\{\fin,\texttt{a},\texttt{l},\texttt{o},\texttt{s}\}$, construya la matriz de conteos $N$ y la matriz de transición $P$ del modelo de bigramas (estimadas por frecuencia relativa).
\end{ejercicio}

\begin{ejercicio}{2}{eje:2-sol-prob}
Con el modelo del \cref{eje:2-sol-matriz}, calcule $P(\text{sol})$, $P(\text{sal})$, $P(\text{los})$ y $P(\text{sos})$. Explique por qué $P(\text{sal})=2P(\text{sol})$.
\end{ejercicio}

\begin{ejercicio}{1}{eje:2-longitud}
En el \cref{ej:bigramas}, la palabra más probable es «a», con $P(\text{a})=0{,}40$. Explique por qué y qué problema plantea para comparar palabras de distinta longitud. Proponga una medida que lo corrija y aplíquela a «a», «ana» y «luna».
\end{ejercicio}

\begin{ejercicio}{3}{eje:2-sin-fin}
Se entrena un modelo de bigramas sin suavizar y sin símbolo de fin, con símbolo de inicio $\langle s\rangle$, sobre las frases «$\langle s\rangle$ a b», «$\langle s\rangle$ b b», «$\langle s\rangle$ b a» y «$\langle s\rangle$ a b». (a) Escriba todas las probabilidades de transición. (b) Calcule la probabilidad de las cuatro frases de dos palabras sobre $\{\texttt{a},\texttt{b}\}$ y compruebe que suman~1. (c) Relacione el resultado con la \cref{prop:longitud-fija}.
\fuente{adaptado de Jurafsky y Martin~\cite{jm3}, cap.~3, ejercicio~3.5 (corpus modificado).}
\end{ejercicio}
